\section{Pinned \texorpdfstring{\code{ConciseIPRA}}{ConciseIPRA} API inventory}

This appendix records the upstream names used by the Lean outline.  It is not a
substitute for the source files; it is a migration checklist.

\begin{longtable}{p{0.29\textwidth}p{0.62\textwidth}}
\toprule
Name & Role \\
\midrule
\endfirsthead
\toprule
Name & Role \\
\midrule
\endhead
\code{PR k} & primitive-recursive function symbols of arity \(k\) \\
\code{PTm n} & terms with \(n\) free variables \\
\code{PTm.app}, \code{PTm.zero}, \code{PTm.succ} & term constructors \\
\code{Fm PTm n} & first-order formulas \\
\code{Fm.QF} & quantifier-free syntactic predicate \\
\code{PRAAx} & iPRA axiom family \\
\code{Prf PRAAx Gamma phi} & object-language derivability \\
\code{PRA.precZero} & primitive recursion zero equation at arbitrary terms \\
\code{PRA.precSucc} & primitive recursion successor equation at arbitrary terms \\
\code{PRA.qfInd} & quantifier-free induction rule \\
\code{Prf.subst} & simultaneous substitution in a derivation \\
\code{Prf.weaken} & context weakening \\
\code{Prf.trans} & simultaneous cut / replacement of hypotheses \\
\code{Prf.impI}, \code{Prf.impE} & implication rules \\
\code{Prf.allI}, \code{Prf.allE} & universal quantifier rules \\
\code{Prf.eqSubst1} & one-variable Leibniz transport \\
\code{Prf.eqSymm}, \code{Prf.eqTrans} & equality symmetry and transitivity \\
\code{Prf.eqCongr}, \code{Prf.eqCongr1} & term congruence \\
\code{NFm.ofPRA} & translation from ordinary iPRA formula to native formula \\
\code{pra_iff_native} & equivalence of ordinary and native derivability \\
\bottomrule
\end{longtable}

The outline intentionally avoids private implementation details of the native
calculus.  It constructs ordinary iPRA proofs first, so future changes to the
compact presentation need affect only the already checked bridge.
